%% ============================================================
%% Oblate Spheroid Excitation Theory (OSET) — Acta Materialia
%% Two-column (elsarticle 5p), full derivations retained
%% pdflatex -> bibtex -> pdflatex -> pdflatex
%% ============================================================
\documentclass[5p,sort&compress]{elsarticle}

\usepackage{amsmath,amssymb,amsfonts}
\usepackage{mathtools}
\usepackage{graphicx}
\usepackage{booktabs}
\usepackage{array}
\usepackage{xcolor}
\usepackage{siunitx}
\usepackage[hidelinks]{hyperref}
\usepackage{microtype}

\graphicspath{{figures/}}
\journal{Acta Materialia}

\newcommand{\OSTZ}{\textsc{ostz}}
\newcommand{\OSET}{\textsc{oset}}
\newcommand{\dF}{\Delta F_0}
\newcommand{\gz}{\gamma_0}
\newcommand{\beff}{b_\mathrm{eff}}
\newcommand{\PN}{Peierls--Nabarro}
\newcommand{\half}{\tfrac{1}{2}}

\begin{document}
\begin{frontmatter}

\title{Oblate Spheroid Excitation Theory (\OSET): A Unified,
  Lattice-Free Foundation for Plastic Deformation from Which
  Dislocations Emerge as Collective Excitations}

\author[iitk]{Albert~Linda\corref{cor}}
\ead{albert.linda@iitk.ac.in}
\author[anna]{K.A.~Padmanabhan}

\cortext[cor]{Corresponding author.}
\affiliation[iitk]{organization={Department of Materials Science and
  Engineering, Indian Institute of Technology Kanpur},
  city={Kanpur}, postcode={208016}, country={India}}
\affiliation[anna]{organization={Professor of Eminence, Materials
  Science and Engineering Programme, Department of Mechanical
  Engineering, Anna University},
  city={Chennai}, postcode={600\,025}, country={India}}

\begin{abstract}
Classical plasticity theory takes the dislocation as its irreducible
elementary entity. This fails for grain boundaries, amorphous solids,
nanocrystals near the glass transition, ceramics, and polymers because
the dislocation is a topological concept that requires a periodic
crystal lattice for its very definition. We propose the
\emph{Oblate Spheroid Excitation Theory} (\OSET), in which the
elementary excitation is a shear-eigenstrained oblate spheroid (the
oblate-spheroidal transformation zone, \OSTZ) embedded in an elastic
continuum. The \OSTZ\ requires no lattice, carries a finite energy
$\dF=\half(\beta_1\gz^2+\beta_2\varepsilon_0^2)GV_0$, and has a
non-singular stress field. Three central results are proved:
(1)~a single \OSTZ\ is equivalent in the far field to a dislocation
dipole; (2)~a coplanar chain of $N$ \OSTZ s is mathematically identical
to a \PN\ dislocation with core width $\zeta=W$ and Burgers vector
$b=N\gz W$; (3)~a full lattice dislocation nucleates when
$N=N_c=b_\mathrm{latt}/(\gz W)$. \OSET\ naturally predicts the
theoretical shear strength, Peierls stress, stacking-fault energy,
core energy, and Frank--Read critical stress, all as derived quantities.
Since dislocations emerge from \OSET\ but \OSET\ does not require
dislocations, \OSET\ is the more fundamental description.
\end{abstract}

\begin{keyword}
oblate spheroid excitation \sep Eshelby inclusion \sep dislocation
emergence \sep grain-boundary sliding \sep Peierls--Nabarro model \sep
superplasticity
\end{keyword}
\end{frontmatter}

%%=========================================================
\section{Introduction}
\label{sec:intro}

\subsection{What dislocation theory requires}

The classical dislocation is defined via the Burgers circuit in a
Bravais lattice:
\begin{equation}
  \oint_C \mathrm{d}u_i = b_i,\qquad b_i\in\Lambda_\mathrm{Bravais}.
  \label{eq:burgers}
\end{equation}
This definition carries three unavoidable restrictions.

\textbf{(R1) Requires a crystal lattice.} No lattice, no Burgers
vector, no dislocation. Grain boundaries, glasses, and polymers are
excluded by definition.

\textbf{(R2) Singular core.} The Volterra stress
$\sigma\sim Gb/2\pi r\to\infty$; a cutoff $r_0\sim b$ is introduced
empirically.

\textbf{(R3) Athermal.} Thermal activation is appended via kink
nucleation; it is not intrinsic.

\subsection{Failure at the nanoscale}

\begin{itemize}
\item \textbf{GBS:} For general high-angle boundaries the DSC lattice
  is commensurate only statistically.
\item \textbf{Nanocrystals at $d<10b$:} Standard plasticity predicts
  hardening not observed (inverse Hall--Petch).
\item \textbf{Metallic glasses:} Deformation via STZs; no dislocations.
\end{itemize}

\subsection{What \OSET\ achieves}

\OSET\ replaces the lattice-dependent, singular, athermal dislocation
with a geometry-defined, non-singular, thermally activated oblate
spheroidal excitation. Dislocations are \emph{derived} as the large-$N$
collective limit. Three inputs: $W$ (OSTZ radius), $\gz$ (shear
eigenstrain), $G$ (shear modulus). All other dislocation quantities
emerge.

%%=========================================================
\section{The Fundamental Entity: the \OSTZ}
\label{sec:ostz}

\subsection{Geometric definition}
\label{sec:geometry}

An \OSTZ\ is an oblate spheroid that has undergone a shear eigenstrain:
\begin{align}
  a_1=a_2=W,\quad a_3=W/2,&\quad \alpha\equiv a_3/a_1=\tfrac{1}{2},
  \notag\\
  V_0=\tfrac{4\pi}{3}a_1^2 a_3&=\tfrac{2\pi}{3}W^3.
  \label{eq:geometry}
\end{align}
The flat face lies in the $x_1$--$x_2$ glide plane. The eigenstrain
for grain-boundary sliding (GBS):
\begin{equation}
  \varepsilon^*_{13}=\varepsilon^*_{31}=\gz/2\;(\text{shear}),\quad
  \varepsilon^*_{ii}=\varepsilon_0/3\;(\text{dilat.})
  \label{eq:eigenstrain}
\end{equation}

This geometry makes $S_{1313}$, not $S_{1212}$, the Eshelby component
governing stored elastic energy.

\subsection{Eshelby solution for the \OSTZ}
\label{sec:eshelby}

By the Eshelby inclusion theorem~\cite{eshelby1957}, the stress inside
a homogeneous ellipsoidal inclusion is \emph{uniform}:
\begin{align}
  \varepsilon_{ij}^\mathrm{in}&=S_{ijkl}\varepsilon^*_{kl},
    \label{eq:eshelby1}\\
  \sigma_{ij}^\mathrm{in}&=C_{ijkl}(S_{klmn}-I_{klmn})\varepsilon^*_{mn}.
    \label{eq:eshelby2}
\end{align}
Uniformity holds only for ellipsoids --- the only shapes for which
$S_{ijkl}$ is spatially uniform.

The relevant Eshelby component for the oblate spheroid~\cite{mura1987}:
\begin{equation}
  S_{1313}=\frac{(1+\alpha^2)I_{13}+(1-2\nu)(I_1+I_3)}{16\pi(1-\nu)},
  \label{eq:S1313}
\end{equation}
where $\Delta(s)=\sqrt{(a_1^2{+}s)(a_2^2{+}s)(a_3^2{+}s)}$,
\begin{align}
  I_\alpha&=2\pi a_1 a_2 a_3\int_0^\infty
    \frac{\mathrm ds}{(a_\alpha^2+s)\Delta(s)},
    \label{eq:Ialpha}\\
  I_{\alpha\beta}&=2\pi a_1 a_2 a_3\int_0^\infty
    \frac{\mathrm ds}{(a_\alpha^2+s)(a_\beta^2+s)\Delta(s)},\notag
\end{align}
with sum rule $I_1+I_2+I_3=4\pi$.

\subsubsection{Derivation of $I_1$}
\label{sec:I1}

\paragraph{Starting point.}
Specialize to $a_1=a_2=1$, $a_3=\alpha$. Then
$\Delta(s)=(1+s)\sqrt{\alpha^2+s}$ and
\begin{equation}
  I_1=2\pi\alpha\int_0^\infty\frac{\mathrm ds}
  {(1+s)^2\sqrt{\alpha^2+s}}.
  \label{eq:I1start}
\end{equation}

\paragraph{Substitution.}
Let $t=\sqrt{\alpha^2+s}$, so $s=t^2-\alpha^2$, $\mathrm ds=2t\,\mathrm dt$,
$\sqrt{\alpha^2+s}=t$, and $1+s=t^2+(1-\alpha^2)\equiv t^2+e^2$
(where $e^2\equiv1-\alpha^2$). New limits: $t:\alpha\to\infty$.
Substituting:
\begin{equation}
  I_1=4\pi\alpha\int_\alpha^\infty\frac{\mathrm dt}{(t^2+e^2)^2}.
  \label{eq:I1sub}
\end{equation}

\paragraph{Antiderivative.}
The reduction formula (proved by differentiating the right side):
\begin{equation}
  \int\frac{\mathrm dt}{(t^2+e^2)^2}=
  \frac{t}{2e^2(t^2+e^2)}+\frac{\arctan(t/e)}{2e^3}+C.
  \label{eq:antideriv}
\end{equation}
\textit{Proof.} Let $F=t/[2e^2(t^2{+}e^2)]$,
$G_f=\arctan(t/e)/(2e^3)$:
\begin{align*}
  F'&=\frac{1}{2e^2}\cdot\frac{(t^2+e^2)-t\cdot2t}{(t^2+e^2)^2}
    =\frac{e^2-t^2}{2e^2(t^2+e^2)^2},\\
  G_f'&=\frac{1}{2e^3}\cdot\frac{1/e}{1+(t/e)^2}
    =\frac{1}{2e^2(t^2+e^2)}.
\end{align*}
Sum: $F'+G_f'=[(e^2-t^2)+(t^2+e^2)]/[2e^2(t^2+e^2)^2]
=1/(t^2+e^2)^2$. $\square$

\paragraph{Evaluating at the bounds.}
\textit{Upper bound} ($t\to\infty$): $F\to0$, $\arctan\to\pi/2$, so
$[F+G_f]_{t\to\infty}=\pi/(4e^3)$.

\textit{Lower bound} ($t=\alpha$): note $\alpha^2+e^2=1$, so
$F|_\alpha=\alpha/(2e^2)$,
$G_f|_\alpha=\arctan(\alpha/e)/(2e^3)$. Hence:
\begin{align}
  \int_\alpha^\infty\frac{\mathrm dt}{(t^2+e^2)^2}
  &=\frac{\pi}{4e^3}-\frac{\alpha}{2e^2}
    -\frac{\arctan(\alpha/e)}{2e^3}\notag\\
  &=\frac{1}{2e^3}\!\left[\frac{\pi}{2}-\arctan\!\left(\frac{\alpha}{e}
    \right)\right]-\frac{\alpha}{2e^2}.
  \label{eq:I1bounds}
\end{align}

\paragraph{Identity $\pi/2-\arctan(\alpha/e)=\arccos\alpha$.}
Set $\alpha=\cos\phi$ ($e=\sin\phi$). Then
$\arctan(\cos\phi/\sin\phi)=\arctan(\cot\phi)=\pi/2-\phi$, so
$\pi/2-\arctan(\alpha/e)=\phi=\arccos\alpha$. Substituting into
Eq.~\eqref{eq:I1bounds} and multiplying by $4\pi\alpha$, and
replacing $e^3=(1-\alpha^2)^{3/2}$:
\begin{equation}
  \boxed{I_1=\frac{2\pi\alpha}{(1-\alpha^2)^{3/2}}
  \!\left[\arccos\alpha-\alpha\sqrt{1-\alpha^2}\right],\quad
  I_3=4\pi-2I_1.}
  \label{eq:I1result}
\end{equation}

\subsubsection{Derivation of $I_{13}=(I_3-I_1)/(1-\alpha^2)$}
\label{sec:I13}

Writing the integrals explicitly ($\Delta=(1+s)\sqrt{\alpha^2+s}$):
\begin{align}
  I_1&=2\pi\alpha\int_0^\infty\frac{\mathrm ds}
    {(1+s)^2\sqrt{\alpha^2+s}},\notag\\
  I_3&=2\pi\alpha\int_0^\infty\frac{\mathrm ds}
    {(1+s)(\alpha^2+s)^{3/2}},\notag\\
  I_{13}&=2\pi\alpha\int_0^\infty\frac{\mathrm ds}
    {(1+s)^2(\alpha^2+s)^{3/2}}.\notag
\end{align}
Forming $I_3-I_1$ and factoring out $1/[(1+s)\sqrt{\alpha^2+s}]$:
\begin{align*}
  I_3-I_1&=2\pi\alpha\int_0^\infty\frac{1}{(1+s)\sqrt{\alpha^2+s}}
  \!\left[\frac{1}{\alpha^2+s}-\frac{1}{1+s}\right]\mathrm ds.
\end{align*}
The bracket equals $(1-\alpha^2)/[(\alpha^2{+}s)(1{+}s)]$:
\begin{align*}
  I_3-I_1&=(1-\alpha^2)\cdot2\pi\alpha\int_0^\infty
  \frac{\mathrm ds}{(1+s)^2(\alpha^2+s)^{3/2}}=(1-\alpha^2)I_{13}.
\end{align*}
Therefore:
\begin{equation}
  \boxed{I_{13}=\frac{I_3-I_1}{1-\alpha^2}.}\quad\square
  \label{eq:I13result}
\end{equation}

\subsubsection{Numerical evaluation ($\alpha=1/2$, $\nu=1/3$)}
\label{sec:numeval}

\textbf{Step~1: compute $I_1$.}
$e=\sqrt{1-0.25}=0.8660$.
Bracket: $\arccos(0.5)-0.5\times0.8660=1.0472-0.4330=0.6142$.
$(0.75)^{3/2}=0.6495$.
$I_1=\pi\times0.6142/0.6495=\mathbf{2.972}$.

\textbf{Step~2: $I_3$ and $I_{13}$.}
$I_3=4\pi-2\times2.972=12.566-5.944=\mathbf{6.622}$.
$I_{13}=(6.622-2.972)/0.75=3.650/0.75=\mathbf{4.867}$.

\textbf{Step~3: $S_{1313}$ from Eq.~\eqref{eq:S1313}.}
Numerator: $1.25\times4.867+(1-2/3)\times9.594
=6.084+3.198=9.282$.
Denominator: $16\pi\times(2/3)=33.51$.
$S_{1313}=9.282/33.51=\mathbf{0.2772}$.

\textbf{Step~4: $\beta_1$.}
$\beta_1=1-2\times0.2772=\mathbf{0.4456}\approx0.446$.

\subsubsection{Limiting-case verification}
\label{sec:limits}

\paragraph{Sphere ($\alpha\to1$).}
For a sphere all semi-axes equal: $I_1=I_2=I_3=4\pi/3$.
Direct integration:
$I_{13}^\text{sphere}=2\pi\int_0^\infty(1+s)^{-7/2}\mathrm ds
=2\pi[u^{-5/2}/(-5/2)]_1^\infty=4\pi/5$.
Substituting ($1+\alpha^2\to2$, $I_1+I_3=8\pi/3$):
\begin{equation}
  S_{1313}^\text{sphere}=\frac{4-5\nu}{15(1-\nu)}.
  \label{eq:S1313sphere}
\end{equation}
For $\nu=1/3$: $S_{1313}^\text{sphere}=7/30=0.233$,
$\beta_1^\text{sphere}=8/15\approx0.533$.

\paragraph{Thin disk ($\alpha\to0$).}
$I_3\to4\pi$, $I_{13}\to4\pi$, $S_{1313}^\text{disk}=\half$,
$\beta_1\to0$: a crack-like disk stores no shear energy.
The OSTZ aspect ratio $\alpha=1/2$ gives the correct intermediate.

\subsubsection{Dilatational constraint factor $\beta_2$}
\label{sec:beta2}

For a spherical inclusion with purely dilatational eigenstrain
$\varepsilon^*_{ij}=(\varepsilon_0/3)\delta_{ij}$, the sphere Eshelby
tensor satisfies:
\begin{equation}
  S_{ijkk}^\text{sphere}=\frac{1+\nu}{3(1-\nu)}\delta_{ij}.
  \label{eq:Sijkk}
\end{equation}
This follows from the known sphere values~\cite{mura1987}
$S_{1111}=(7-5\nu)/[15(1-\nu)]$,
$S_{1122}=(5\nu-1)/[15(1-\nu)]$, summing:
$S_{1111}+S_{1122}+S_{1133}=(1+\nu)/[3(1-\nu)]$. Full derivation
of $\beta_2$ follows in §\ref{sec:Edilat}; the result:
\begin{equation}
  \boxed{\beta_2=\frac{4(1+\nu)}{9(1-\nu)}.}
  \label{eq:beta2}
\end{equation}
For $\nu=1/3$: $\beta_2=4(4/3)/[9(2/3)]=0.889$.

Two formulas for $\beta_2$ exist: (a) exact oblate formula
$\beta_2^\text{oblate}=1-S_{3333}(1-2\nu)/[2(1-\nu)]=0.816$;
(b) spherical dilatation approximation $\beta_2^\text{sphere}=0.889$~\cite{pad-part1}.
They differ by $\sim9\%$ at $\nu=1/3$;
since $\beta_2\varepsilon_0^2\ll\beta_1\gz^2$, the difference is
$<2\%$ of $\dF$ and inconsequential.

\begin{table}[htbp]
\centering
\small
\caption{Eshelby components for $\alpha=1/2$.}
\label{tab:eshelby}
\begin{tabular}{ccccc}
\toprule
$\nu$ & $S_{1313}$ & $S_{1212}$ & $S_{3333}$ & $\beta_1$\\
\midrule
0.30 & 0.2822 & 0.1769 & 0.7264 & 0.4356\\
1/3  & 0.2772 & 0.1739 & 0.7364 & 0.4456\\
0.35 & 0.2745 & 0.1723 & 0.7418 & 0.4510\\
0.40 & 0.2656 & 0.1670 & 0.7596 & 0.4688\\
\bottomrule
\end{tabular}
\end{table}

\subsection{Note on $\beta_1$: two physical regimes}

\textbf{(a) Eshelby (elastic) value:}
$\beta_1^\text{Esh}=1-2S_{1313}\approx0.44$--$0.49$
for $\alpha=0.5$, $\nu=0.28$--$0.44$.

\textbf{(b) Padmanabhan et al.\ effective value:} $\beta_1^\text{eff}\approx1$.

\textit{Physical reconciliation.} OSTZs at grain boundaries are embedded
in GB material with $G_\text{GB}\approx G/2$ (GB softening~\cite{wolf1990}):
$\beta_1^\text{eff}=\beta_1^\text{Esh}\times G/G_\text{GB}
\approx0.45\times2=0.90\approx1$.

\subsection{OSTZ elastic strain energy}
\label{sec:energy}

The Eshelby energy theorem:
$\dF=-\half\sigma_{ij}^\mathrm{in}\varepsilon^*_{ij}V_0$.

\subsubsection{Part~1: Shear contribution}
\label{sec:Eshear}

Eigenstrain $\varepsilon^*_{13}=\varepsilon^*_{31}=\gz/2$.

\textit{Step~1 (constrained strain).}
By minor symmetry $S_{1313}=S_{1331}$:
$\varepsilon_{13}^\mathrm{in}=S_{1313}(\gz/2+\gz/2)=S_{1313}\gz$.

\textit{Step~2 (elastic strain).}
$e_{13}^\mathrm{el}=\varepsilon_{13}^\mathrm{in}-\varepsilon^*_{13}
=(S_{1313}-\half)\gz$.

\textit{Step~3 (interior stress).}
$\sigma_{13}^\mathrm{in}=2Ge_{13}^\mathrm{el}
=2G(S_{1313}-\half)\gz$.

\textit{Step~4 (energy contraction).}
Both $(1,3)$ and $(3,1)$ contribute equally:
\[
  \dF^\text{shear}=-\half(\sigma_{13}^\mathrm{in}\varepsilon^*_{13}
  +\sigma_{31}^\mathrm{in}\varepsilon^*_{31})V_0
  =-\half\sigma_{13}^\mathrm{in}\gz V_0.
\]
Substituting:
$\dF^\text{shear}=-\half\cdot2G(S_{1313}-\half)\gz\cdot\gz V_0
=G(\half-S_{1313})\gz^2 V_0$.
\begin{equation}
  \boxed{\dF^\text{shear}=\half\beta_1 G\gz^2 V_0,\quad
  \beta_1\equiv1-2S_{1313}.}
  \label{eq:Eshear}
\end{equation}
Since $S_{1313}<\half$ for any physical oblate spheroid,
$\beta_1>0$ and the shear energy is positive definite.

\subsubsection{Part~2: Dilatational contribution}
\label{sec:Edilat}

Eigenstrain $\varepsilon^*_{ij}=(\varepsilon_0/3)\delta_{ij}$.

\textit{Step~1 (constrained strain inside a sphere).}
From Eqs.~\eqref{eq:eshelby1} and \eqref{eq:Sijkk}:
\begin{equation}
  \varepsilon_{ij}^\mathrm{in}=S_{ijkl}\varepsilon^*_{kl}
  =\frac{\varepsilon_0}{3}\cdot\frac{1+\nu}{3(1-\nu)}\delta_{ij}
  =\frac{(1+\nu)\varepsilon_0}{9(1-\nu)}\delta_{ij}.
  \label{eq:constrained_dil}
\end{equation}

\textit{Step~2 (elastic strain).}
\begin{align}
  e_{ij}^\mathrm{el}&=\varepsilon_{ij}^\mathrm{in}-\varepsilon^*_{ij}
    =\varepsilon_0\!\left[\frac{1+\nu}{9(1-\nu)}-\frac{1}{3}\right]\delta_{ij}
    \notag\\
  &=\varepsilon_0\cdot\frac{(1+\nu)-3(1-\nu)}{9(1-\nu)}\delta_{ij}
    =-\frac{2(1-2\nu)\varepsilon_0}{9(1-\nu)}\delta_{ij}.
    \label{eq:edilelastic}
\end{align}
(Numerator: $1+\nu-3+3\nu=4\nu-2=-2(1-2\nu)$.)
Volumetric trace:
$e_{kk}^\mathrm{el}=-6(1-2\nu)\varepsilon_0/[9(1-\nu)]
=-2(1-2\nu)\varepsilon_0/[3(1-\nu)]$.

\textit{Step~3 (interior stress).}
Using $\lambda=2G\nu/(1-2\nu)$:
\begin{align}
  \lambda e_{kk}^\mathrm{el}
    &=\frac{2G\nu}{1-2\nu}\cdot\frac{-2(1-2\nu)\varepsilon_0}{3(1-\nu)}
    =-\frac{4G\nu\varepsilon_0}{3(1-\nu)},\notag\\
  2Ge_{11}^\mathrm{el}
    &=-\frac{4G(1-2\nu)\varepsilon_0}{9(1-\nu)}.
    \notag
\end{align}
Adding with common denominator $9(1-\nu)$:
\begin{align}
  \sigma_{11}^\mathrm{in}&=-\frac{12G\nu\varepsilon_0}{9(1-\nu)}
    -\frac{4G(1-2\nu)\varepsilon_0}{9(1-\nu)}\notag\\
  &=-\frac{G\varepsilon_0}{9(1-\nu)}[12\nu+4-8\nu]
    =-\frac{4G(1+\nu)\varepsilon_0}{9(1-\nu)}.
    \label{eq:sig11in}
\end{align}

\textit{Step~4 (energy contraction).}
Three diagonal components contribute equally:
$\sigma_{ij}^\mathrm{in}\varepsilon^*_{ij}
=3\sigma_{11}^\mathrm{in}(\varepsilon_0/3)=\sigma_{11}^\mathrm{in}\varepsilon_0$.
\begin{equation}
  \boxed{\dF^\text{dilat}=\half\beta_2 G\varepsilon_0^2 V_0,\quad
  \beta_2=\frac{4(1+\nu)}{9(1-\nu)}.}
  \label{eq:Edilat}
\end{equation}

\subsubsection{Total activation energy}

Shear and dilatational eigenstrains involve orthogonal tensor
components ($\varepsilon^*_{13}$ vs.\ $\varepsilon^*_{ii}$);
their cross-terms vanish upon contraction with the isotropic stiffness
tensor. Therefore:
\begin{equation}
  \boxed{\dF=\half(\beta_1\gz^2+\beta_2\varepsilon_0^2)GV_0,\quad
  V_0=\tfrac{2\pi}{3}W^3.}
  \label{eq:dF0}
\end{equation}
\textit{Numerical check:}
$\dF=\half(0.00446+0.00222)\times2.2\times10^{10}\times0.884\times10^{-27}
=6.5\times10^{-20}$~J$=0.406$~eV$\approx0.38$~eV.

\begin{table}[htbp]
\centering
\small
\caption{Canonical \OSTZ\ parameters.}
\label{tab:canonical}
\begin{tabular}{lccc}
\toprule
Parameter & Symbol & GB/superpl. & Crystal\\
\midrule
OSTZ radius & $W$ & $2.5b$ & $\approx b$\\
Shear eigenstrain & $\gz$ & 0.10 & 0.12\\
Dilat.\ eigenstrain & $\varepsilon_0$ & 0.05 & ---\\
Shear constraint & $\beta_1$ & 0.446 & $\approx1$\\
Dilat.\ constraint & $\beta_2$ & 0.889 & ---\\
Activation energy & $\dF$ & 0.38~eV & ---\\
Critical number & $N_c$ & 4 & $\approx8$\\
\bottomrule
\end{tabular}
\end{table}

%%=========================================================
\section{Far-Field Stress: The Dipole Structure}
\label{sec:dipole}

\subsection{Why a dipole, not a monopole}

The force monopole vanishes:
\begin{equation}
  F_k=\oint_{\partial V_0}\sigma_{ij}^* n_j\,\mathrm dS
  =C_{ijkl}\varepsilon^*_{kl}
  \underbrace{\oint_{\partial V_0}n_j\,\mathrm dS}_{=0}=0.
  \label{eq:monopole}
\end{equation}
(The surface integral of the outward normal over any closed surface
vanishes by the divergence theorem.) The leading far-field term is
therefore the \emph{force-dipole} at order $r^{-3}$.

The full multipole expansion~\cite{mura1987} is
\begin{multline}
  \sigma_{ij}(\mathbf r)=
  \underbrace{F_k\mathcal G_{ijk}}_{\text{monopole }r^{-2}}
  +\underbrace{M_{kl}\partial_l\mathcal G_{ijk}}_{\text{dipole }r^{-3}}
  \\+\underbrace{Q_{klm}\partial_l\partial_m\mathcal G_{ijk}}_{r^{-4}}
  +\cdots
  \label{eq:multipole}
\end{multline}
Since $F_k=0$, the leading nonzero term is the dipole. The moment
tensor~\cite{mura1987} is
\begin{equation}
  M_{kl}=C_{ijkl}\varepsilon^*_{ij}V_0\;\Rightarrow\;
  M_{13}=M_{31}=G\gz V_0,
  \label{eq:moment}
\end{equation}
describing a double couple: two force pairs of equal and opposite sign,
separated by the OSTZ diameter $2W$.

\subsection{Step A: Kelvin Green's function and its derivative}
\label{sec:stepA}

The elastic displacement due to a unit point force at the origin in an
infinite isotropic medium (Kelvin--Somigliana solution~\cite{mura1987}):
\begin{equation}
  \mathcal G_{ij}(\mathbf r)=A\!\left[\frac{(3-4\nu)\delta_{ij}}{r}
  +\frac{x_ix_j}{r^3}\right],\quad A=\frac{1}{16\pi G(1-\nu)}.
  \label{eq:Gij}
\end{equation}

\paragraph{Differentiating the first term.}
Since $\partial_l(1/r)=-x_l/r^3$:
\[
  \partial_l\!\left[\frac{(3-4\nu)\delta_{ij}}{r}\right]
  =-\frac{(3-4\nu)\delta_{ij}x_l}{r^3}.
\]

\paragraph{Differentiating the second term (product rule).}
$\partial x_i/\partial x_l=\delta_{il}$, $\partial(r^{-3})/\partial x_l=-3x_l/r^5$:
\[
  \partial_l\!\left(\frac{x_ix_j}{r^3}\right)
  =\frac{\delta_{il}x_j+\delta_{jl}x_i}{r^3}-\frac{3x_lx_ix_j}{r^5}.
\]

\paragraph{Combined result.}
\begin{align}
  \frac{\partial\mathcal G_{ij}}{\partial x_l}
  =A\Bigl[&-\frac{(3-4\nu)\delta_{ij}x_l}{r^3}
    +\frac{\delta_{il}x_j+\delta_{jl}x_i}{r^3}
    \notag\\
  &-\frac{3x_lx_ix_j}{r^5}\Bigr].
  \label{eq:dG}
\end{align}

\subsection{Step B: Displacement from the moment tensor}
\label{sec:stepB}

For an Eshelby eigenstrain inclusion, the far-field exterior
displacement is~\cite{eshelby1957,mura1987} $u_i=-M_{kl}\partial_l\mathcal G_{ik}$.
With $M_{13}=M_{31}=G\gz V_0$ (all other components zero),
the only nonzero contributions are $(k,l)=(1,3)$ and $(3,1)$.

Substituting Eq.~\eqref{eq:dG} (with $\delta_{13}=0$):
\begin{align}
  \frac{\partial\mathcal G_{i1}}{\partial x_3}
  &=A\!\left[-\frac{(3-4\nu)\delta_{i1}x_3}{r^3}
    +\frac{\delta_{i3}x_1}{r^3}-\frac{3x_3x_ix_1}{r^5}\right],
    \notag\\
  \frac{\partial\mathcal G_{i3}}{\partial x_1}
  &=A\!\left[-\frac{(3-4\nu)\delta_{i3}x_1}{r^3}
    +\frac{\delta_{i1}x_3}{r^3}-\frac{3x_1x_ix_3}{r^5}\right].
    \notag
\end{align}

Adding term by term, the $\delta_{i1}x_3$ coefficient is
$1-(3-4\nu)=4\nu-2=2(2\nu-1)$ (same for $\delta_{i3}x_1$),
and the $r^{-5}$ terms add to give $-6Ax_1x_3x_i/r^5$:
\begin{multline}
  u_i=-G\gz V_0 A\Bigl[\frac{2(2\nu-1)(\delta_{i1}x_3+\delta_{i3}x_1)}{r^3}
  \\-\frac{6x_1x_3x_i}{r^5}\Bigr].
  \label{eq:ui}
\end{multline}

Components relevant to $\sigma_{13}$:
\begin{align}
  u_1&=-G\gz V_0 A\!\left[\frac{2(2\nu-1)x_3}{r^3}
    -\frac{6x_1^2x_3}{r^5}\right],
    \label{eq:u1}\\
  u_3&=-G\gz V_0 A\!\left[\frac{2(2\nu-1)x_1}{r^3}
    -\frac{6x_1x_3^2}{r^5}\right].
    \label{eq:u3}
\end{align}

\subsection{Step C: Stress from displacement gradients}
\label{sec:stepC}

For the shear component ($\delta_{13}=0$, so the $\lambda$-term in
Hooke's law vanishes): $\sigma_{13}=G(\partial_3u_1+\partial_1u_3)$.

\paragraph{Auxiliary results.}
$\partial_3(x_3/r^3)=(r^2-3x_3^2)/r^5$;
$\partial_3(x_3/r^5)=(r^2-5x_3^2)/r^7$; likewise for $x_1$.

\paragraph{Computing $\partial_3u_1$.}
Differentiating Eq.~\eqref{eq:u1} with respect to $x_3$:
\begin{align}
  \partial_3u_1&=-G\gz V_0 A\Bigl[\frac{2(2\nu-1)(r^2-3x_3^2)}{r^5}
    \notag\\
  &\phantom{=-G\gz V_0 A\Bigl[}\;-\frac{6x_1^2(r^2-5x_3^2)}{r^7}\Bigr].
    \label{eq:du1dx3}
\end{align}

\paragraph{Computing $\partial_1u_3$.}
By the same procedure (swap $x_1\leftrightarrow x_3$):
\begin{align}
  \partial_1u_3&=-G\gz V_0 A\Bigl[\frac{2(2\nu-1)(r^2-3x_1^2)}{r^5}
    \notag\\
  &\phantom{=-G\gz V_0 A\Bigl[}\;-\frac{6x_3^2(r^2-5x_1^2)}{r^7}\Bigr].
    \label{eq:du3dx1}
\end{align}

\paragraph{Summing and converting to unit vectors $\hat r_i=x_i/r$.}
$1/r^5$ group:
$2(2\nu-1)[2-3(\hat r_1^2+\hat r_3^2)]/r^3$.
$1/r^7$ group, expanding $(x_1^2+x_3^2)r^2-10x_1^2x_3^2$:
$[-6(\hat r_1^2+\hat r_3^2)+60\hat r_1^2\hat r_3^2]/r^3$.
Collecting the $(\hat r_1^2+\hat r_3^2)$ terms:
$-6(2\nu-1)-6=-12\nu$ (verify: $12\nu-6+6=12\nu$). Hence:
\begin{multline}
  \partial_3u_1+\partial_1u_3=
  -G\gz V_0 A\cdot\frac{1}{r^3}\\
  \times\bigl[4(2\nu-1)-12\nu(\hat r_1^2+\hat r_3^2)
  +60\hat r_1^2\hat r_3^2\bigr].
  \label{eq:gradsum}
\end{multline}

\subsection{Step D: Angular tensor and far-field formula}
\label{sec:stepD}

Multiplying Eq.~\eqref{eq:gradsum} by $G$ and inserting
$A=1/[16\pi G(1-\nu)]$:
\[
  \sigma_{13}=\frac{G\gz V_0}{16\pi(1-\nu)r^3}
  \bigl[-4(2\nu-1)+12\nu(\hat r_1^2+\hat r_3^2)-60\hat r_1^2\hat r_3^2
  \bigr].
\]
Writing this as $G\gz V_0\mathcal T_{13}/[2\pi(1-\nu)r^3]$ requires
dividing the bracket by 8 (ratio of prefactors $2\pi/16\pi=1/8$):
\begin{equation}
  \boxed{\mathcal T_{13}(\hat{\mathbf r})
  =\tfrac{1}{2}-\nu+\tfrac{3\nu}{2}(\hat r_1^2+\hat r_3^2)
  -\tfrac{15}{2}\hat r_1^2\hat r_3^2.}
  \label{eq:Tij}
\end{equation}
\begin{equation}
  \boxed{\sigma_{13}(\mathbf r)=\frac{G\gz V_0}{2\pi(1-\nu)}
  \cdot\frac{\mathcal T_{13}(\hat{\mathbf r})}{r^3}.}
  \label{eq:dipole}
\end{equation}
The $r^{-3}$ decay follows from differentiating the Kelvin $r^{-1}$
Green's function twice: two orders faster than the Volterra $r^{-1}$
field.

\subsection{Step E: Verification in special directions}
\label{sec:stepE}

\paragraph{Glide plane ($x_3=0$, $\hat r_3=0$).}
$\mathcal T_{13}|_{\hat r_3=0}=\half-\nu+(3\nu/2)\hat r_1^2$.
Expanding $r^2=x_1^2+x_2^2$ and collecting:
\begin{equation}
  \sigma_{13}(x_1,x_2,0)=\frac{G\gz V_0}{4\pi(1-\nu)r^5}
  [(1+\nu)x_1^2+(1-2\nu)x_2^2].
  \label{eq:glidecheck}
\end{equation}
Positive definite for $0<\nu<\half$. $\checkmark$

\paragraph{On-axis ($\hat r_1=0$, $\hat r_3=1$).}
$\mathcal T_{13}=\half-\nu+3\nu/2=(1+\nu)/2>0$. By oblate axial
symmetry, the same result holds along $\hat x_1$. $\checkmark$

\paragraph{Nodal surface.}
Setting $\mathcal T_{13}=0$ yields no real solution for
$0<\nu<\half$: there is no nodal surface in the half-space
$x_3\ge0$. The \OSTZ\ promotes shear in all glide-plane directions
(cooperative nucleus).

\subsection{Glide-plane stress: six-step derivation}
\label{sec:glideplane}

Working at $x_3=0$ ($r^2=x^2+y^2$, $x\equiv x_1$, $y\equiv x_2$).

\textbf{Step~1 (off-diagonal $\mathcal G_{13}$).}
$\partial^2(\mathcal G_{13})/(\partial x_1\partial x_3)$ at $z=0$:
gives $(y^2-2x^2)/r^5$ (using $\partial_3(x_1x_3/r^3)|_{z=0}
=x_1/r^3$, then $\partial_1$).

\textbf{Step~2 (diagonal $\mathcal G_{33}$).}
At $z=0$: $\mathcal G_{33}|_{z=0}=A(3-4\nu)/r$.
$\partial_1^2(1/r)=(3x^2-r^2)/r^5=(2x^2-y^2)/r^5$:
\[
  \partial_1^2\mathcal G_{33}\big|_{z=0}=A(3-4\nu)(2x^2-y^2)/r^5.
\]

\textbf{Step~3 (diagonal $\mathcal G_{11}$).}
At $z=0$: $\mathcal G_{11}|_{z=0}=A[(3-4\nu)/r+x^2/r^3]$.
$\partial_3^2(1/r)|_{z=0}=-1/r^3$;
$\partial_3^2(1/r^3)|_{z=0}=-3/r^5$:
\[
  \partial_3^2\mathcal G_{11}\big|_{z=0}=A[-(3-4\nu)/r^3-3x^2/r^5].
\]
Expanding: $-(3-4\nu)r^2-3x^2=-(6-4\nu)x^2-(3-4\nu)y^2$.

\textbf{Step~4 (collect numerator, units of $A/r^5$).}
\begin{align*}
  x^2\text{ coeff.:}&\;2(y^2\!-\!2x^2)\to-4\\
  &+(6\!-\!8\nu)-(3\!-\!4\nu)-3=-4(1+\nu),\\
  y^2\text{ coeff.:}&\;2-(3\!-\!4\nu)-(3\!-\!4\nu)=-4(1-2\nu).
\end{align*}
Numerator: $-4A[(1+\nu)x^2+(1-2\nu)y^2]/r^5$.

\textbf{Step~5 (assemble).}
\[
  \sigma_{13}(x,y,0)=\frac{G\gz V_0[(1+\nu)x^2+(1-2\nu)y^2]}
  {4\pi(1-\nu)r^5}.
\]

\textbf{Step~6 (re-express with $\beff$, restore finite $W$).}
$G\gz V_0=G(\gz)\tfrac{2\pi W^3}{3}=G\beff\tfrac{2\pi W^2}{3}$
(using $\beff=\gz W$). Then $G\gz V_0/[4\pi(1-\nu)]=G\beff W^2/[6(1-\nu)]$.
Regularising $r^2\to r^2+W^2$:
\begin{equation}
  \boxed{\sigma_{13}(x,y,0)=
  \frac{G\beff W^2[(1+\nu)x^2+(1-2\nu)y^2]}
  {6(1-\nu)(x^2+y^2+W^2)^{5/2}}.}
  \label{eq:glideplane}
\end{equation}
Positive definite; finite at the origin.

\paragraph{Physical interpretation.}
Along $\hat x$ ($y=0$): promotes forward shear.
Along $\hat y$ ($x=0$): promotes lateral shear.
No nodal planes in the glide plane: the \OSTZ\ is a
\emph{cooperative nucleus}.

\subsection{Effective Burgers vector}
\label{sec:beff}

\textbf{Method~1 (kinematic, exact):}
$\beff=\gz\times2a_3=\gz W$.

\textbf{Method~2 (moment tensor):}
$M_{13}=G\gz V_0=G\beff\pi W^2$
$\Rightarrow\beff=2\gz W/3$ (underestimates 33\%).

\textbf{Method~3 (Lorentzian):}
$\beff$ fixed by normalisation; not an independent check.

Primary definition: $\boxed{\beff=\gz W}$.
\label{eq:beff}

\begin{table}[htbp]
\centering
\small
\caption{Single \OSTZ\ vs.\ Volterra dislocation.}
\label{tab:compare1}
\begin{tabular}{p{1.6cm}p{2.3cm}p{2.3cm}}
\toprule
Property & Single \OSTZ & Volterra\\
\midrule
Stress decay & $r^{-3}$ (dipole) & $r^{-1}$ (monopole)\\
Topo.\ charge & 0 (neutral) & $\pm b$\\
Net $b$ & 0 & $b$\\
Core stress & Finite ($\le G\gz$) & $\to\infty$\\
Influence & $\sim W$ & $\infty$\\
Elastic $E$ & $\half\beta_1\gz^2GV_0$ & $Gb^2L\ln(R/r_0)$\\
Thermal act. & Intrinsic & Appended\\
\bottomrule
\end{tabular}
\end{table}

%%=========================================================
\section{Emergence of Dislocations from \OSTZ\ Chains}
\label{sec:emergence}

\subsection{Lorentzian dislocation density: four steps}
\label{sec:lorentzian}

\subsubsection{Step A: 2D dislocation density (BCS compliance)}
\label{sec:stepA-BCS}

The Bilby--Cottrell--Swinden (BCS) compliance~\cite{bilby1955}
$\delta u_1=2(1-\nu)\sigma_{13}/G$ relates local slip to traction.
The shear stress in the isotropic approximation:
\begin{equation}
  \sigma_{13}(x,y,0)\approx
  \frac{G\beff W^2}{2\pi(1-\nu)(x^2+y^2+W^2)^{3/2}}.
  \label{eq:gpisotropic}
\end{equation}
(This replaces the exact angular form by a spherically symmetric
approximation, valid because the true field $\sim r^{-3}$ converges
absolutely and is positive everywhere in the glide plane.)
Carrying out the cancellations
($G$ with $1/G$; $2(1-\nu)$ with $2\pi(1-\nu)$ leaving $1/\pi$):
\begin{equation}
  \rho_{2D}(x,y)=\frac{\beff W^2}{\pi(x^2+y^2+W^2)^{3/2}}.
  \label{eq:rho2D}
\end{equation}

\textit{Normalisation check.}
Converting to polar coordinates $r^2=x^2+y^2$:
\begin{align*}
  \iint\rho_{2D}\,\mathrm dx\,\mathrm dy
  &=\frac{\beff W^2}{\pi}\cdot2\pi
    \int_0^\infty\frac{r\,\mathrm dr}{(r^2+W^2)^{3/2}}.
\end{align*}
Substitution $u=r^2+W^2$, $du=2r\,dr$:
$\int=\half[-2/\sqrt{u}]_{W^2}^\infty=1/W$.
Therefore $\iint\rho_{2D}=2\beff W$.
(The result is $2\beff W$, not $\beff$, because $\rho_{2D}$ is slip
per unit glide-plane area; the Abel projection recovers the correct
normalisation below.)

\subsubsection{Step B: Abel projection}
\label{sec:stepB-abel}

Integrate over $y$ with $a^2=x^2+W^2$; let $y=a\tan\theta$:
\begin{equation}
  \int_{-\infty}^\infty\frac{\mathrm dy}{(a^2+y^2)^{3/2}}
  =\frac{1}{a^2}\!\int_{-\pi/2}^{\pi/2}\cos\theta\,\mathrm d\theta
  =\frac{2}{x^2+W^2}.
  \label{eq:abel}
\end{equation}
Raw 1D density: $\rho_1^\text{raw}=2\beff W^2/[\pi(x^2+W^2)]$,
which integrates to $2\beff W$.

\subsubsection{Step C: Lorentzian (normalised)}
\label{sec:stepC-lorentz}

Divide by $2W$ to enforce $\int\rho_1\,\mathrm dx=\beff$:
\begin{equation}
  \boxed{\rho_1(x)=\frac{\beff}{\pi}\cdot\frac{W}{x^2+W^2}.}
  \label{eq:rho1}
\end{equation}
\textit{Check:}
$\int_{-\infty}^\infty\rho_1\,\mathrm dx=(\beff W/\pi)(\pi/W)=\beff$.
$\square$

The Lorentzian is the unique non-negative distribution integrating to
$\beff$ with width $W$ and reducing to a Dirac delta as $W\to0$.
Both \OSET\ and the P--N model are governed by the same
Cauchy--Hilbert Green's operator of planar elasticity.

\subsubsection{Step D: Fourier transform}
\label{sec:stepD-fourier}

For $k>0$, close the contour in the lower half-plane (where
$e^{-ikx}\to0$). The only enclosed pole is at $x=-iW$:
\[
  \mathrm{Res}\!\left[\frac{We^{-ikx}}{\pi(x^2+W^2)},\,-iW\right]
  =\frac{e^{-kW}}{-2iW}.
\]
Clockwise contour ($-2\pi i\times\text{Res}$):
\begin{equation}
  \hat\rho(k)=\beff\,e^{-|k|W}.
  \label{eq:rhohat}
\end{equation}
At $k=2\pi/b$: $\hat\rho=\beff e^{-2\pi W/b}$ (Peierls lattice-filter
factor). The OSTZ core acts as a low-pass filter: high-spatial-frequency
lattice components are exponentially suppressed.

\subsection{Theorem: the $N$-\OSTZ\ chain is a \PN\ dislocation}
\label{sec:PN}

For $N$ co-planar \OSTZ s:
$\rho_N(x)=N\cdot(\beff/\pi)W/(x^2+W^2)$.

The cumulative slip from $-\infty$ to $x$:
\begin{align}
  U(x)&=\int_{-\infty}^x\rho_N(x')\,\mathrm dx'\notag\\
  &=\frac{N\beff W}{\pi}\cdot\frac{1}{W}
    \!\left[\arctan\frac{x'}{W}\right]_{-\infty}^x
    \notag\\
  &=\frac{N\beff}{2}+\frac{N\beff}{\pi}\arctan\frac{x}{W}.
  \label{eq:slip}
\end{align}
(Using $\arctan(-\infty)=-\pi/2$.)

Setting $b\equiv N\beff$ and $\zeta\equiv W$:
\begin{equation}
  U(x)=\frac{b}{2}+\frac{b}{\pi}\arctan\frac{x}{\zeta}.
  \label{eq:PNprofile}
\end{equation}
This is exactly the \PN\ profile~\cite{peierls1940,nabarro1947}.
\textit{Limits:} $U(-\infty)=0$, $U(+\infty)=b$, $U(0)=b/2$. $\square$

The continuum dislocation density is $\rho(x)=\partial U/\partial x$:
\[
  \rho(x)=\frac{N\beff}{\pi}\cdot\frac{W}{x^2+W^2}=\rho_N(x).\quad\square
\]

\textbf{Corollary.} The core half-width $\zeta=W$ is derived from
\OSTZ\ geometry, not a free parameter. The \PN\ model is not an
independent theory: it is the continuum limit of the \OSTZ\ ensemble.

\subsection{Chain stress field and Volterra limit}
\label{sec:chainStress}

\begin{equation}
  \sigma_{13}^{(N)}(x)=\frac{G}{2\pi(1-\nu)}\,\mathrm{P.V.}
  \int_{-\infty}^\infty\frac{\rho_N(x')}{x-x'}\,\mathrm dx'.
  \label{eq:conv}
\end{equation}

\paragraph{Partial fractions.}
Decompose $W/[(x'^2+W^2)(x-x')]$:
$W=A(x'^2+W^2)+(Bx'+C)(x-x')$.
Matching powers of $x'$:
$x'^2$: $A=B$;
$x'^1$: $Bx-C=0\Rightarrow C=Ax$;
constant: $A(x^2+W^2)=W\Rightarrow A=W/(x^2+W^2)$.
Hence:
\[
  \frac{W}{(x'^2+W^2)(x-x')}
  =\frac{W}{x^2+W^2}\!\left[\frac{1}{x-x'}
  +\frac{x'+x}{x'^2+W^2}\right].
\]

\paragraph{Term~1.}
P.V.\,$\int\mathrm dx'/(x-x')=0$ (integrand is odd about $x'=x$).

\paragraph{Term~2.}
$\int x'/(x'^2+W^2)\,\mathrm dx'=0$ (odd in $x'$).
$\int\mathrm dx'/(x'^2+W^2)=\pi/W$
(via $x'=W\tan\theta$).
So Term~2 gives $\pi x/(x^2+W^2)$.

\paragraph{Combining:}
\begin{equation}
  \boxed{\sigma_{13}^{(N)}(x)=\frac{GN\beff}{2\pi(1-\nu)}
  \cdot\frac{x}{x^2+W^2}.}
  \label{eq:chainStress}
\end{equation}
Non-singular at $x=0$; maximum at $x=W$.

\paragraph{Volterra limit ($W\to0$).}
$\lim_{W\to0}x/(x^2+W^2)=1/x$, so:
\begin{equation}
  \sigma_{13}^\text{Volterra}=\frac{Gb}{2\pi(1-\nu)x}.
  \label{eq:volterra}
\end{equation}
The Volterra singularity is a collective limit; no single \OSTZ\ is
singular.

\paragraph{Burgers circuit.}
$\oint_C\mathrm du_1=N\beff=b$. Quantisation
$b\in\Lambda_\mathrm{Bravais}$ arises when the lattice constrains $N$
to $N_c$.

\subsection{Critical cluster number}
\label{sec:Nc}

Setting $N\beff=b_\mathrm{latt}$:
\begin{equation}
  \boxed{N_c=\frac{b_\mathrm{latt}}{\gz W}.}
  \label{eq:Nc}
\end{equation}
FCC Cu ($W=b$, $\gz=0.12$): $N_c\approx8$.
GB regime ($W=2.5b$, $\gz=0.10$): $N_c=4$
(confirmed experimentally in Zn--22Al~\cite{astanin-part2}).

Why dislocations do not form in amorphous materials: no periodic
potential to lock displacements at quantised values
$b_i\in\Lambda_\mathrm{Bravais}$, so OSTZs activate individually
or in small clusters ($N\ll N_c$): STZ-like deformation.

%%=========================================================
\section{Dislocation Quantities as Derived Results}
\label{sec:derived}

\subsection{Theoretical shear strength}
\label{sec:strength}

The cascade instability threshold: $\tau^*\gz V_0=\dF=\half\beta_1\gz^2GV_0$:
\begin{equation}
  \boxed{\tau^*=\frac{\beta_1\gz G}{2}.}
  \label{eq:strength}
\end{equation}
For $\beta_1\approx1$, $\gz=0.12$: $\tau^*\approx G/17$
(Frenkel range $G/10$--$G/30$).

\subsection{Peierls stress}
\label{sec:peierls}

\subsubsection{OSET-derived amplitude}

Setting $\tau^*=(2\pi/b)V_0^\text{OSET}$:
\begin{equation}
  V_0^\text{OSET}=\frac{\beta_1\gz Gb}{4\pi}.
  \label{eq:V0OSET}
\end{equation}
For Cu: $\tau^*\approx G/17$ vs.\ the Frenkel estimate $G/[2(1-\nu)]\approx0.75G$.

\subsubsection{Fourier convolution}

Misfit energy $E_\text{mis}=\int V(x-x_0)\rho(x)\,\mathrm dx$,
$V=V_0[1-\cos(2\pi x/b)]$.
Using $\hat\rho(k_1)=be^{-2\pi W/b}$ at $k_1=2\pi/b$
and Euler's formula $\cos=\mathrm{Re}[e^{i\theta}]$:
\begin{equation}
  E_\text{mis}(x_0)\big|_\text{osc}=
  -\frac{Gb^2}{4\pi(1-\nu)}\,e^{-2\pi W/b}
  \cos\!\left(\frac{2\pi x_0}{b}\right).
  \label{eq:Emisfit}
\end{equation}
(Constant $V_0 b$ term dropped as it plays no role in $\tau_P$.)

\subsubsection{Peierls barrier and half-space correction}

Barrier: $\Delta E_P=2Ae^{-2\pi W/b}$ with $A=Gb^2/[4\pi(1-\nu)]$.
Maximum slope: $\tau_P=(2\pi A/b^2)e^{-2\pi W/b}=G e^{-2\pi W/b}/[2(1-\nu)]$.
With half-space correction~\cite{peierls1940} ($W\to W/(1-\nu)$, prefactor $\times2$):
\begin{equation}
  \boxed{\tau_P=\frac{2G}{1-\nu}
  \exp\!\left(-\frac{2\pi W}{b(1-\nu)}\right).}
  \label{eq:peierls}
\end{equation}
Cu ($W=b$, $\nu=1/3$): $\tau_P\approx2.1\times10^{-4}G$
(experiment: $\sim10^{-4}G$).

\subsection{Core energy}
\label{sec:coreE}

$E_\text{core}=N_c\dF$. Substituting $N_c=b/(\gz W)$ and
$\dF=\pi\beta_1\gz^2GW^3/3$:
\begin{equation}
  E_\text{core}=\frac{b}{\gz W}\cdot\frac{\pi\beta_1\gz^2GW^3}{3}
  =\frac{\pi\beta_1 b\gz^{\cancel2} GW^{\cancel3}}
  {3\cancel\gz\cancel{W}}.
  \label{eq:coreEderiv}
\end{equation}
\begin{equation}
  \boxed{E_\text{core}=\frac{\pi\beta_1 b\gz GW^2}{3}.}
  \label{eq:coreE}
\end{equation}
Writing $E_\text{core}=\alpha Gb^2$: for Cu ($W\approx b$),
$\alpha=\pi\beta_1\gz/3\approx0.13$ (exp/DFT: $0.1$--$0.2$).

\subsection{Frank--Read critical stress}
\label{sec:FR}

Following Frank \& Read~\cite{frank-read1950}:

\textbf{Step~1 (line tension).}
$T=Gb^2/[4\pi(1-\nu)]$. Semicircle of radius $R=L/2$ has arc
length $\pi L/2$:
\begin{equation}
  E_\text{bow}=T\cdot\pi R=\frac{Gb^2L}{8(1-\nu)}.
  \label{eq:Ebow}
\end{equation}

\textbf{Step~2 (external work).}
Enclosed area $\pi R^2/2=\pi L^2/8$:
\begin{equation}
  W_\text{ext}=\tau_\text{FR}\cdot b\cdot\frac{\pi L^2}{8}.
  \label{eq:Wext}
\end{equation}

\textbf{Step~3 ($W_\text{ext}=E_\text{bow}$).}
Dividing both sides by $bL/8$:
$\tau_\text{FR}\cdot\pi L=Gb/(1-\nu)$:
\begin{equation}
  \boxed{\tau_\text{FR}=\frac{Gb}{\pi(1-\nu)L}\approx\frac{Gb}{2L}.}
  \label{eq:FR}
\end{equation}
Cu at $L=1\;\mu$m: 14~MPa vs.\ classical 12.4~MPa (13\%).

\subsection{Stacking-fault energy}
\label{sec:SFE}

\subsubsection{FCC Burgers ratios}

$|b_\text{full}|=a_0/\sqrt2$; $|b_\text{partial}|=a_0/\sqrt6$:
\begin{equation}
  \frac{|b_\text{partial}|}{|b_\text{full}|}
  =\frac{\sqrt2}{\sqrt6}=\frac{1}{\sqrt3}.
  \label{eq:bratio}
\end{equation}

\subsubsection{\OSTZ\ cluster for the Shockley partial}

Setting $N_p\gz W=|b_\text{partial}|=N_c\gz W/\sqrt3$:
$N_p=N_c/\sqrt3$.

Evaluating $(N_c-N_p)/N_p$:
\[
  N_c-N_p=N_c\!\left(1-\frac{1}{\sqrt3}\right),\quad
  \frac{N_c-N_p}{N_p}=\sqrt3\!\left(1-\frac{1}{\sqrt3}\right)=\sqrt3-1.
\]

\subsubsection{SFE derivation}

The energy per fault area ($\pi W^2$ per OSTZ site, $N_p$ sites):
\[
  \gamma_\text{SF}=\frac{(N_c-N_p)\dF}{\pi W^2 N_p}
  =\frac{(\sqrt3-1)\dF}{\pi W^2}.
\]
Substituting $\dF=\pi\beta_1\gz^2GW^3/3$ and cancelling
$W^3/W^2=W$:
\begin{equation}
  \boxed{\gamma_\text{SF}=\frac{(\sqrt3-1)\beta_1\gz^2GW}{3}.}
  \label{eq:SFE}
\end{equation}
Cu ($G=48.3$~GPa, $W=b=0.2556$~nm, $\gz=0.12$, $\beta_1=1$):
$\gamma_\text{SF}=43$~mJ/m$^2$ vs.\ experiment $45$~mJ/m$^2$ (5\%).
Equation~\eqref{eq:SFE} identifies $W\approx b$ as the reason for
$\gamma_\text{SF}\propto Gb$ across FCC metals.

\begin{table}[htbp]
\centering
\small
\caption{SFE predictions ($\gz=0.12$, $W=b$, $\beta_1^\text{eff}=1$).}
\label{tab:crystal}
\begin{tabular}{lcccc}
\toprule
Metal & $G$ (GPa) & $b$ (nm) & $\gamma_\text{SF}^\text{OSET}$ & Exp.\\
      &  & & (mJ/m$^2$) & (mJ/m$^2$)\\
\midrule
Cu & 48.3 & 0.256 & 43 & 45\\
Al & 26.2 & 0.286 & 26 & 166\\
Ni & 76.0 & 0.249 & 67 & 125\\
Ag & 30.3 & 0.289 & 31 & 16\\
Au & 27.0 & 0.288 & 28 & 32\\
\bottomrule
\end{tabular}
\end{table}

%%=========================================================
\section{Grain-Boundary Energy}
\label{sec:gb}

\subsection{Coherent boundaries (quadratic in eigenstrain)}

A coherent twin boundary (CTB) carries a coherency strain over an
interfacial area. Its energy is \emph{quadratic} in the eigenstrain
--- identical to the \OSET\ SFE~\eqref{eq:SFE}:
$\gamma_\text{CTB}=(\sqrt3-1)\beta_1\gz^2GW/3$.

\subsection{Incoherent boundaries: the Read--Shockley law}
\label{sec:RS}

A low/high-angle boundary is a dislocation array spaced $D=b/\theta$.
In \OSET\ each dislocation is an $N_c$-\OSTZ\ chain; summing the elastic
dipole fields reproduces the classical dislocation line energy
$E_d=(Gb^2/4\pi(1-\nu))\ln(R/r_0)$. Dividing by $D$ with
$R\sim D/2$ ($\ln(R/r_0)\to A-\ln\theta$) yields the
Read--Shockley law~\cite{readshockley1950}:
\begin{align}
  \gamma_\text{GB}(\theta)&=E_0\,\theta(A-\ln\theta),
    \label{eq:RS}\\
  E_0&=\frac{Gb}{4\pi(1-\nu)},\quad A=1+\ln\theta_m.\notag
\end{align}
The GB energy scales linearly with $b$ (dislocation array = linear in
eigenstrain via $b=N_c\gz W$), not quadratically with $\gz$
(activation barrier). Conflating these objects causes 5--40$\times$
error.

%%=========================================================
\section{\OSTZ\ Hamiltonian and Statistical Mechanics}
\label{sec:stat}

\subsection{Construction of the Hamiltonian}
\label{sec:Hamiltonian}

The grain boundary is modelled as a 2D array of sites $i$, each with
binary occupation $n_i\in\{0,1\}$.

\paragraph{Term~1: Self-energy.}
\begin{equation}
  \mathcal H_\text{self}=\sum_i\dF_i\,n_i.
  \label{eq:Hself}
\end{equation}

\paragraph{Term~2: Mechanical work.}
Applied stress $\tau$ does work $(\tau-\tau_0)\gz V_0$ per activated
site, lowering the barrier ($\tau_0\propto L^{-1/2}$ is the threshold
back-stress):
\begin{equation}
  \mathcal H_\text{mech}=-\sum_i(\tau-\tau_0)\gz V_0\,n_i.
  \label{eq:Hmech}
\end{equation}

\paragraph{Term~3: Elastic interaction.}
Two activated \OSTZ s at $\mathbf x_i$ and $\mathbf x_j$ interact via
their overlapping $r^{-3}$ fields:
\begin{equation}
  E_{ij}^\text{int}=-\sigma_{13}^{(i)}(\mathbf x_j)\gz V_0
  =-\frac{G\gz^2V_0^2\mathcal T_{13}(\hat{\mathbf r}_{ij})}
  {2\pi(1-\nu)r_{ij}^3}.
  \label{eq:Eint}
\end{equation}
The leading (dipole--dipole) term defines the interaction kernel:
\begin{align}
  J_{ij}&=G\gz^2V_0^2\,\mathcal K\!\left(\frac{r_{ij}}{W}\right),
    \label{eq:Jij}\\
  \mathcal K(\rho)&=\frac{3\cos^2\theta_{ij}-1}{4\pi(1-\nu)\rho^3}.
    \notag
\end{align}
$J_{ij}>0$ (cooperative) for \OSTZ s aligned along the slip direction
($\theta_{ij}=0$); $J_{ij}<0$ (anti-cooperative) perpendicular.
The form is identical to the magnetic dipole--dipole interaction: a
direct consequence of the shared $r^{-3}$ Green's function structure.
\begin{equation}
  \mathcal H_\text{int}=-\half\sum_{i\ne j}J_{ij}\,n_in_j.
  \label{eq:Hint}
\end{equation}

\paragraph{Full Hamiltonian.}
\begin{multline}
  \mathcal H=\sum_i\!\left[\dF_i-(\tau-\tau_0)\gz V_0\right]n_i\\
  -\half\sum_{i\ne j}J_{ij}\,n_in_j.
  \label{eq:Hamiltonian}
\end{multline}
This is formally a \emph{random-field Ising model} with quenched
disorder $\dF_i$ and ferromagnetic exchange $J_{ij}$.

\subsection{Mean-field decoupling}
\label{sec:MF}

\paragraph{The linearisation step.}
Write $n_i=\bar n+\delta n_i$. Expand $n_in_j$ and drop
$\delta n_i\delta n_j$ (Bragg--Williams approximation~\cite{bragg1934}):
\begin{equation}
  n_in_j\approx n_i\bar n+n_j\bar n-\bar n^2.
  \label{eq:BW}
\end{equation}

\paragraph{Effect on the interaction term.}
Substituting into $-\half\sum_{i\ne j}J_{ij}n_in_j$; the first two
terms are equal by relabelling ($i\leftrightarrow j$); with
translation invariance $\sum_{j\ne i}J_{ij}=zJ_0$:
\begin{equation}
  -\half\sum_{i\ne j}J_{ij}n_in_j\approx
  -z\bar nJ_0\sum_i n_i+\text{const.}
  \label{eq:MFinteraction}
\end{equation}

\paragraph{Mean-field Hamiltonian.}
The total factorises into independent single-site problems:
\begin{align}
  \mathcal H^\text{MF}&=\sum_i h^\text{eff}n_i+\text{const},
    \label{eq:heff}\\
  h^\text{eff}&=\dF-(\tau-\tau_0)\gz V_0-z\bar nJ_0.\notag
\end{align}
The effective barrier $h^\text{eff}$ is the bare cost $\dF$ reduced
by the mechanical driving force and the cooperative elastic
interaction from already-activated neighbours.

\subsubsection{Self-consistency equation}
\label{sec:selfconsistency}

Single-site partition function $\mathcal Z_i=1+e^{-h^\text{eff}/kT}$.
Mean occupation:
\begin{equation}
  \bar n=\langle n_i\rangle
  =\frac{e^{-h^\text{eff}/kT}}{1+e^{-h^\text{eff}/kT}}.
  \label{eq:selfconsistency}
\end{equation}
Substituting $h^\text{eff}$ explicitly:
\begin{equation}
  \boxed{\bar n=\frac{\exp\!\left[-\tfrac{\dF-(\tau-\tau_0)\gz V_0
    -z\bar nJ_0}{kT}\right]}{1+\exp\!\left[-\tfrac{\dF-(\tau-\tau_0)
    \gz V_0-z\bar nJ_0}{kT}\right]}}
  \label{eq:selfconsfull}
\end{equation}
This nonlinear equation is solved iteratively for $\bar n$.

In the \emph{dilute limit} ($h^\text{eff}\gg kT$, $z\bar nJ_0\to0$),
$\bar n\approx\exp(-\dF/kT)\cdot\exp((\tau-\tau_0)\gz V_0/kT)$
--- only the forward activation. The full forward--backward balance
is captured by the Eyring kinetic route (§\ref{sec:sinh}).

\subsection{Sinh flow law: step-by-step derivation}
\label{sec:sinh}

\paragraph{Step~1 (strain per activation).}
Each \OSTZ\ activation shifts the material by $\beff=\gz W$ over a
cross-section $\pi W^2$ in a grain of size $d$. The macroscopic
strain increment:
\begin{equation}
  \delta\gz_\text{macro}=\frac{\beff\cdot\pi W^2}{\pi W^2\cdot d}
  =\frac{\gz W}{d}.
  \label{eq:deltaGamma}
\end{equation}

\paragraph{Step~2 (Eyring forward/backward rates).}
Following Eyring~\cite{eyring1935}, bias $\Delta w=(\tau-\tau_0)\gz V_0/2$.
Forward barrier $\dF-\Delta w$; reverse $\dF+\Delta w$:
\begin{equation}
  \Gamma^\pm=\nu_0\exp\!\left(-\frac{\dF\mp\Delta w}{kT}\right).
  \label{eq:Gamma}
\end{equation}

\paragraph{Step~3 (factor out $e^{-\dF/kT}$).}
Using $e^{a+b}=e^a\cdot e^b$:
\begin{align*}
  e^{-(\dF-\Delta w)/kT}&=e^{-\dF/kT}\cdot e^{+\Delta w/kT},\\
  e^{-(\dF+\Delta w)/kT}&=e^{-\dF/kT}\cdot e^{-\Delta w/kT}.
\end{align*}

\paragraph{Step~4 ($e^x-e^{-x}=2\sinh x$).}
\begin{multline}
  \dot N_\text{net}=\Gamma^+-\Gamma^-\\
  =2\nu_0\sinh\!\left[\frac{(\tau-\tau_0)\gz V_0}{2kT}\right]
  e^{-\dF/kT}.
  \label{eq:Nnet}
\end{multline}

\paragraph{Step~5 (macroscopic rate).}
Multiply by $\delta\gz_\text{macro}=\gz W/d$:
\begin{equation}
  \boxed{\dot\gz=\frac{2W\gz\nu_0}{d}
  \sinh\!\left[\frac{(\tau-\tau_0)\gz V_0}{2kT}\right]
  e^{-\dF/kT}.}
  \label{eq:rate}
\end{equation}
Identical to Padmanabhan et al.\ (1996, Eq.~8)~\cite{pad-part1}.

\subsubsection{Comparison: partition-function route}

From $\ln\mathcal Z=N_\text{sites}\ln(1+e^{-h^\text{eff}/kT})$,
differentiation with respect to $(\tau-\tau_0)\gz V_0/kT$ gives
$N_\text{sites}\bar n$. In the dilute limit, this is a
\emph{forward-bias-only} expression --- not equal to $2e^{-\dF/kT}
\sinh(\Delta w/kT)$.
The sinh rate equation is an Eyring \emph{kinetic} result.
Detailed balance: $\Gamma^+/\Gamma^-=e^{2\Delta w/kT}$. $\checkmark$

\subsubsection{Connection to Padmanabhan et al.\ (1996)}

With $\Delta w=(\tau-\tau_0)\gz V_0/2$ and $(2W\gz/d)\dot w=\dot\gz$,
Eq.~(8) of~\cite{pad-part1} is recovered exactly.
Including a lognormal disorder distribution $f(\dF)$:
\begin{multline}
  \dot\gz=\frac{2W\gz\nu_0}{d}\int_0^\infty f(\dF)\\
  \times\sinh\!\left[\frac{(\tau-\tau_0)\gz V_0}{2kT}\right]
  e^{-\dF/kT}\,\mathrm d(\dF).
  \label{eq:rateDisorder}
\end{multline}
Universal lognormal parameter $\bar A=5.6021$~\cite{venkatesh-part4};
CV$=14.8\%$; range $[0.28,0.52]$~eV.

\subsection{\OSTZ\ interaction as the origin of Taylor hardening}
\label{sec:taylor}

When $z\bar nJ_0$ is non-negligible, the effective threshold
stress increases:
\begin{equation}
  \tau_0^\text{eff}=\tau_0+\frac{z\bar nJ_0}{\gz V_0}.
  \label{eq:tau0eff}
\end{equation}
With $\bar n=\pi W^2\rho_\text{disl}$ and
$J_0\approx G\gz^2V_0^2/[4\pi(1-\nu)W^3]$,
substituting $V_0=(2\pi/3)W^3$:
\begin{equation}
  \tau_0^\text{eff}-\tau_0=
  \frac{z\gz G\pi W^2}{6(1-\nu)}\,\rho_\text{disl}.
  \label{eq:Taylorhardening}
\end{equation}
\OSET\ predicts hardening \emph{linear} in $\rho_\text{disl}$
(stage-I regime). The classical Taylor $\sqrt\rho$ law~\cite{taylor1934} requires
the additional geometric forest-obstacle argument
($\tau_c=Gb/L$, $L\propto\rho^{-1/2}$) and is not derivable
from the \OSTZ\ Ising Hamiltonian alone.

%%=========================================================
\section{Why \OSET\ is More Fundamental}
\label{sec:fundamental}

\subsection{Logical priority theorem}

\textbf{Theorem.} \textit{\OSET$\Rightarrow$dislocation theory,
but not conversely.}

($\Rightarrow$): Sections~\ref{sec:PN}--\ref{sec:chainStress}
derive the Volterra field with quantised $b=N_c\gz W$.

($\not\Rightarrow$): Dislocation theory cannot describe
(a)~single \OSTZ\ events;
(b)~\OSTZ\ excitations in amorphous matter;
(c)~GBS without a DSC lattice;
(d)~the nanocrystal-to-glass transition. $\square$

\subsection{Ontological hierarchy}

\begin{center}
\small
$\text{OSTZ}
\xrightarrow{N=1}\text{STZ/GBS}
\xrightarrow{N=N_p}\text{partial disln}$\\[3pt]
$\xrightarrow{N=N_c}\text{full disln}
\xrightarrow{N\gg N_c}\text{slip band}$
\end{center}

Dislocations occupy one branch of the \OSTZ\ tree. \OSET\ is the root.

\begin{table}[htbp]
\centering
\small
\caption{Classical dislocation theory vs.\ \OSET.}
\label{tab:comparison}
\begin{tabular}{p{1.5cm}p{2.2cm}p{2.3cm}}
\toprule
Property & Classical & \OSET\\
\midrule
Entity & Topological line & Eshelby spheroid\\
Lattice & Required & Not required\\
Materials & Crystals & All\\
Core stress & $\to\infty$ & $\le G\gz$\\
Core energy & Fitted $\alpha Gb^2$ & Eq.~\eqref{eq:coreE}\\
Core width & Fitted & $\zeta=W$\\
Burgers $b$ & Input & $N_c\gz W$\\
Thermal act. & Appended & Intrinsic\\
$\tau_P$ & $W$ fitted & Eq.~\eqref{eq:peierls}\\
$\tau^*$ & Frenkel sep. & Eq.~\eqref{eq:strength}\\
SFE & Fitted & Eq.~\eqref{eq:SFE}\\
GB plast. & DSC needed & Natural\\
\bottomrule
\end{tabular}
\end{table}

%%=========================================================
\section{\OSET\ Across Material Classes}
\label{sec:classes}

\textbf{Perfect crystal ($\dF\gg kT$).}
OSTZs are rare; act in clusters of $N_c$; recover Orowan equation.

\textbf{Grain boundary / superplastic.}
OSTZs at GB free-volume sites; $N\ll N_c$; rate equation =
Padmanabhan et al.~\cite{pad-part1} Eq.~(39) directly.

\textbf{Metallic glass (no $N_c$).}
$N=1$ STZ picture~\cite{argon1979,falk-langer1998}. At low $T$:
localised deformation (shear bands = OSTZ percolation paths).
At high $T$: homogeneous flow.

\textbf{Nanocrystal near the glass transition ($d\to d_0$).}
As $d\to d_0=2\sqrt6 W$, the threshold $\tau_0\to0$ and
$N_c^\text{eff}(d)\to1$: glass-like behaviour. \OSET\ naturally
spans the nanocrystal-to-glass transition.

%%=========================================================
\section{Novel Predictions}
\label{sec:predictions}

\textbf{P1: Core width $\zeta=W$} (testable by HAADF-STEM imaging of
dislocation cores).

\textbf{P2: Core energy without fitting.}
$E_\text{core}/\text{length}=\pi\beta_1 b\gz GW/6$;
Cu: $0.124$~eV/\AA\ (DFT: $0.05$--$0.15$~eV/\AA).

\textbf{P3: SFE scales as $\gz^2GW\approx\gz^2Gb$}, explaining
$\gamma_\text{SF}\propto Gb$ across FCC metals.

\textbf{P4: Linear hardening}
$\tau_0^\text{eff}\propto\rho_\text{disl}$ (Eq.~\eqref{eq:Taylorhardening})
at low $T$; decreases with $T$ through $\gz(T)$ and $J_0\propto G(T)$.

\textbf{P5: Sub-dislocation displacement bursts} $\beff\approx0.1b$,
accessible with picometer-precision AFM-based nanoindentation.

%%=========================================================
\section{Unified Rate Equation}
\label{sec:unified}

Combining the \OSTZ\ partition function (§\ref{sec:stat}) with
disorder-field theory (activation energy variance $\sigma_F^2$)
and crystalline crossover $\Theta$:
\begin{multline}
  \dot\gz=\frac{2W\bar\gz\nu_0}{d}
  \sinh\!\left[\frac{(\tau-\tau_0^*)\bar\gz V_0}{2kT}\right]\\
  \times\exp\!\!\left(-\frac{\Delta\bar F_0}{kT}
  +\frac{\sigma_F^2}{2(kT)^2}\right)\Theta(N_c,T,d),
  \label{eq:rateUnified}
\end{multline}
where $\tau_0^*=\tau_0\sqrt{1+\rho_{G\gz}\sigma_G\sigma_\gz}$ and
\begin{equation}
  \Theta=\begin{cases}
    1 & \text{glass/GB regime}\\
    N_c^{-1}\!\sum_{N=1}^{N_c}\!N\,e^{-E_c(N)/kT} & \text{crystal}
  \end{cases}
  \label{eq:Theta}
\end{equation}
In the crystal limit $\Theta\to N_c$ (Orowan); in the glass/GB limit
$\Theta\to1$ (Padmanabhan et al.~\cite{pad-part1}, Eq.~(39)).

%%=========================================================
\section{Validation}
\label{sec:validation}

\subsection{Validation against Padmanabhan et al.\ (1996)}

\textbf{T1} (Eshelby tensor): $\beta_1=0.446$, $\beta_2=0.889$. Exact.

\textbf{T2} (Canonical): $\dF=0.405$ vs.\ $0.38$~eV (6\%).

\textbf{T3} ($W=2.5b$):
$\delta V_\text{free}=\varepsilon_0 V_0\approx0.1\Omega$/GB atom;
confirms Wolf MD~\cite{wolf1990}.

\textbf{T4} ($N_c=4$): $b/(\gz W)=0.3/(0.1\times0.75)=4.0$ exactly;
confirmed by experiment in Zn--22Al~\cite{astanin-part2}.

\textbf{T5} (Rate equation): Al-12Si at 773~K,
$d=10\;\mu$m, $\tau-\tau_0=5$~MPa:
$\Delta w=2.2\times10^{-22}$~J$=0.00138$~eV;
$\sinh(0.0207)=0.0207$; $e^{-5.71}=0.00330$;
$\dot\gz=0.16$~s$^{-1}$ (superplastic window $10^{-3}$--$10^{-1}$~s$^{-1}$).

\textbf{T6} (Threshold stress scaling $\tau_0\propto L^{-1/2}$):
Al-33Cu ratio predicted 2.0, measured 1.8--2.1.

\textbf{T7} (Lognormal): $\bar A=5.6021$, CV$=14.8\%$,
range $[0.28,0.52]$~eV.

\begin{table}[htbp]
\centering
\small
\caption{Predictions validated against
  Padmanabhan et al.\ (1996).}
\label{tab:1996validation}
\begin{tabular}{p{2.2cm}cc}
\toprule
Claim & Pred. & Agr.\\
\midrule
$\beta_1=0.446$ & $1-2\times0.2772$ & Exact\\
$\beta_2=0.889$ & $(4/9)(4/3)/(2/3)$ & Exact\\
$\gz=0.10$ & Wolf MD, bubble-raft & Exact\\
$\varepsilon_0=0.05$ & $\delta V\approx0.1\Omega$/atom & Exact\\
$N_c=4$ & $4.0$ derived & Exact (expt.)\\
$\dF=0.38$~eV & 0.405~eV & 6\%\\
Rate eq.\ sinh & Eyring kinetics & Exact\\
$\tau_0\propto L^{-1/2}$ & ratio 2.0 & $\le10\%$\\
$Q\sim80\text{--}120$~kJ/mol & 90--120 & $\le30\%$\\
\bottomrule
\end{tabular}
\end{table}

\subsection{Post-1996 literature}

\textbf{L1} Sripathi \& Padmanabhan~\cite{sripathi2014}: 33 systems,
$R=0.89$--$1.00$.

\textbf{L2} $\beta_1$ convention reconciled:
$0.446\times0.01=0.00446\approx1.779\times0.0025=0.00445$.

\textbf{L3} $\gz=0.10$: Wolf MD~\cite{wolf1990} ($0.08$--$0.12$), bubble-raft,
Argon~\cite{argon1979} ($0.08$--$0.14$), Falk--Langer~\cite{falk-langer1998} ($0.10$).
Four independent methods.

\textbf{L4} BMGs~\cite{buenz2015}: 8 systems,
$\dot\varepsilon_\text{pred}/\dot\varepsilon_\text{exp}=1.1$--$2.5$.

\textbf{L5} Nano-Ni~\cite{pad2018}:
$8\times10^{-4}$ vs.\ $5\times10^{-4}$--$2\times10^{-3}$~s$^{-1}$
(factor 2).

\textbf{L6} High-strain-rate superplasticity~\cite{padbasariya2009}: 8 alloy/composite
systems; all within $2\times$.

\textbf{L9} Peierls--Nabarro~\cite{nabarro1997}: Fe exact;
Cu within $2\times$.

\textbf{L10} Lorentzian~\cite{vermeulen1995}:
Al at $\rho=10^{13}$~m$^{-2}$, $M=0.16<0.5$ (Lorentzian regime).

\subsection{Harisankar--Padmanabhan (2025): 41 systems}

\begin{table}[htbp]
\centering
\small
\caption{\OSET\ vs.\ Harisankar--Padmanabhan~\cite{harisankar2025}.}
\label{tab:hp}
\begin{tabular}{lccc}
\toprule
Quantity & \OSET & Paper (mean) & Agr.\\
\midrule
$\varepsilon_0$ & 0.05 & $0.049\pm0.012$ & $<2\%$\\
$\gz$ & 0.10--0.12 & $0.085\pm0.020$ & 15\%\\
$\beta_1\gz^2+\beta_2\varepsilon_0^2$ & 0.0067 & 0.0057 & 15\%\\
$Q/\dF$ & 1 & $\approx4$ & see text\\
$\gamma_B$ (J/m$^2$) & $E_0\theta_m$ & 0.30--1.70 & $4.5\times$\\
\bottomrule
\end{tabular}
\end{table}

\textbf{Key findings:}
(1)~$\varepsilon_0$ predicted to $<2\%$;
(2)~$\gz$ mean 0.085 lies between the GB ($0.10$) and crystal ($0.12$)
canonical values, consistent with temperature softening
$\gz(T)=0.0827+1.342/T$;
(3)~$\dF\approx Q/4$: the remaining three-quarters is
non-elastic (GB migration, triple-junction accommodation,
cooperative rearrangement of $N_a\approx16$ boundary units);
(4)~$\gamma_B$ to median $4.5\times$ --- a constant multiplicative
offset showing no correlation with homologous temperature
($r=0.00$), confirming correct $Gb$ scaling.

%%=========================================================
\section{Discussion}
\label{sec:discussion}

\OSET\ establishes the ontological hierarchy: single \OSTZ\ = elastic
dipole = GBS/STZ event; $N$-chain = \PN\ dislocation;
$N=N_c$ = lattice dislocation.

\paragraph{Strengths.}
FCC noble/transition metals (Cu, Ni, Au, Ag): SFE ($5\%$ for Cu),
Peierls stress (within $2\times$), core energy --- all without free
parameters. GB/superplastic rate equation validated across
41 material systems.

\paragraph{Limitations.}
(i)~Al: anomalously high SFE requires $\gz\approx0.30$ (NFE band
structure).
(ii)~BCC metals: Peierls stress underpredicted $10$--$100\times$
(BCC core has $W<b$; using $W=0.2b$ recovers experiment).
(iii)~$Q/\dF\approx4$: explained by mesoscopic $N_a\approx16$
($Q=N_a\dF$) plus lognormal tail convolution; the elastic \OSTZ\
barrier is 25\% of the total.

%%=========================================================
\section{Conclusions}
\label{sec:conclusions}

\begin{enumerate}

\item \OSTZ\ is lattice-free, non-singular, thermally activated.
  Eshelby parameters $\beta_1=0.446$, $\beta_2=0.889$ are derived
  rigorously from $I_1$, $I_{13}$
  (Sections~\ref{sec:I1}--\ref{sec:numeval}).

\item Dislocations emerge as \OSTZ\ collectives:
  dipole Eq.~\eqref{eq:dipole}, \PN\ profile Eq.~\eqref{eq:PNprofile},
  non-singular chain stress Eq.~\eqref{eq:chainStress}; Volterra
  singularity only as $W\to0$, Eq.~\eqref{eq:volterra}.

\item Five dislocation observables are derived without adjustable
  parameters: $\tau^*$~\eqref{eq:strength}, $\tau_P$~\eqref{eq:peierls},
  $E_\text{core}$~\eqref{eq:coreE}, $\tau_\text{FR}$~\eqref{eq:FR},
  $\gamma_\text{SF}$~\eqref{eq:SFE}.

\item The \OSTZ\ Hamiltonian~\eqref{eq:Hamiltonian} yields the sinh
  flow law~\eqref{eq:rate} and linear
  hardening~\eqref{eq:Taylorhardening}.

\item Validation: $\varepsilon_0$ within $2\%$, $\gz$ within $15\%$,
  $\gamma_B$ to median $4.5\times$, across 41 superplastic systems.

\item \OSET$\Rightarrow$dislocation theory (not conversely):
  \OSET\ is the more fundamental description
  (Table~\ref{tab:comparison}).

\end{enumerate}

\section*{CRediT}
\textbf{A.~Linda:} Methodology, Formal analysis, Writing.
\textbf{K.A.~Padmanabhan:} Conceptualization, Supervision, Review.

\section*{Competing interests}
None.

\section*{Data availability}
Jupyter notebook with all derivations and numerical verifications
accompanies this article.

\begin{thebibliography}{99}

\bibitem{eshelby1957}
J.D. Eshelby, Proc.\ R. Soc.\ Lond.\ A \textbf{241} (1957) 376.

\bibitem{mura1987}
T. Mura, \textit{Micromechanics of Defects in Solids},
2nd ed., Kluwer, 1987.

\bibitem{peierls1940}
R.E. Peierls, Proc.\ Phys.\ Soc.\ \textbf{52} (1940) 34.

\bibitem{nabarro1947}
F.R.N. Nabarro, Proc.\ Phys.\ Soc.\ \textbf{59} (1947) 256.

\bibitem{argon1979}
A.S. Argon, Acta Metall.\ \textbf{27} (1979) 47.

\bibitem{falk-langer1998}
M.L. Falk, J.S. Langer, Phys.\ Rev.\ E \textbf{57} (1998) 7192.

\bibitem{pad-part1}
K.A. Padmanabhan, H.J. Maier, H. Gleiter,
Mater.\ Sci.\ Technol.\ \textbf{12} (1996) 391.

\bibitem{astanin-part2}
V.S. Astanin, S.N. Faizova, K.A. Padmanabhan,
Mater.\ Sci.\ Technol.\ \textbf{12} (1996) 489.

\bibitem{venkatesh-part4}
B. Venkatesh, K.A. Padmanabhan et al.,
Mater.\ Sci.\ Technol.\ \textbf{12} (1996) 545.

\bibitem{sripathi2014}
S. Sripathi, K.A. Padmanabhan,
J.\ Mater.\ Sci.\ \textbf{49} (2014) 2085.

\bibitem{buenz2015}
J. Buenz, K.A. Padmanabhan, G. Wilde,
Intermetallics \textbf{60} (2015) 7.

\bibitem{harisankar2025}
K.R. Harisankar, K.A. Padmanabhan,
Mater.\ Sci.\ Eng.\ A \textbf{930} (2025) 148175.

\bibitem{readshockley1950}
W.T. Read, W. Shockley, Phys.\ Rev.\ \textbf{78} (1950) 275.

\bibitem{frank-read1950}
F.C. Frank, W.T. Read, Phys.\ Rev.\ \textbf{79} (1950) 722.

\bibitem{nabarro1997}
F.R.N. Nabarro, Philos.\ Mag.\ A \textbf{75} (1997) 703.

\bibitem{pad2018}
K.A. Padmanabhan, Mater.\ Sci.\ Eng.\ A (2018).

\bibitem{wolf1990}
D. Wolf, Acta Metall.\ Mater.\ \textbf{38} (1990) 781.

\bibitem{vermeulen1995}
A.C. Vermeulen et al., J.\ Appl.\ Phys.\ \textbf{77} (1995) 5026.

\bibitem{woodward2002}
C. Woodward, S.I. Rao,
Phys.\ Rev.\ Lett.\ \textbf{88} (2002) 216402.

\bibitem{padgleiter2014}
K.A. Padmanabhan, H. Gleiter,
Beilstein J.\ Nanotechnol.\ \textbf{5} (2014) 1502.

\bibitem{bilby1955}
B.A. Bilby, A.H. Cottrell, K.H. Swinden,
Proc.\ R. Soc.\ Lond.\ A \textbf{272} (1963) 304.

\bibitem{bragg1934}
W.L. Bragg, E.J. Williams,
Proc.\ R. Soc.\ Lond.\ A \textbf{145} (1934) 699.

\bibitem{eyring1935}
H. Eyring, J. Chem.\ Phys.\ \textbf{3} (1935) 107.

\bibitem{taylor1934}
G.I. Taylor,
Proc.\ R. Soc.\ Lond.\ A \textbf{145} (1934) 362.

\bibitem{padbasariya2009}
K.A. Padmanabhan, M.R. Basariya,
Mater.\ Sci.\ Eng.\ A \textbf{527} (2009) 225.

\end{thebibliography}

\end{document}
